\section{Introduction}


Discussions of consciousness typically begin with questions such as: where is it generated, which circuits are responsible, and which representational format does it correspond to? This paper changes the framing of these questions. Consciousness is not treated here as something produced somewhere; it is modeled as a state that opens when processes at different time scales become mutually referenceable while maintaining a certain delay and coupling.

In this view, consciousness is not identical to the self, to thought, or to judgment. It is the condition under which fast predictive processes become referenceable to the layers of slow memory, reporting, and sedimentation. Accordingly, the aim of this paper is not to provide a final definition of consciousness, but to establish a structural position from which consciousness can be rigorously discussed.

Previous supplementary treatments addressed qualia, Libet experiments, split-brain phenomena, measurement, and the localization problem individually. This paper reorganizes them into a single independent theoretical draft.

\subsection{The Position of This Paper: Three Declarations}


This section states in advance the position this paper takes with respect to existing theories of consciousness, restricted to three points. Detailed comparison is given in §18. The purpose here is to fix the coordinate frame that the reader should maintain throughout the paper.

\textbf{First: an opening model, not a generative model}

This paper does not treat consciousness as an output generated somewhere. Many models — including Global Workspace, IIT, and Higher-Order theories — describe consciousness as the product of a process or as the endpoint of an integration measure. This paper does not adopt that premise as its starting point.

More precisely, consciousness is not modeled here as a generated object. It is treated as a state (a log-referencing state) that opens when a specific speed difference, coupling viscosity, internal delay, and meaning gradient are established between fast predictive processes and slow sedimentary memory processes. This difference is not merely definitional; it concerns the structure of the question itself — what should be asked about consciousness.

\textbf{Second: description as a temporal regime, not reduction to a single quantity}

This paper does not reduce consciousness to a single physical quantity, information measure, or integration index. Consciousness is not reducible to any one of: speed difference $\Delta v$, coupling viscosity $\mu$, internal delay $\tau$, effective torque-like quantity $\mathcal{T}_{\Phi}$, or meaning gradient $\nabla\Phi$. Consciousness is a temporal non-closure regime that obtains when these variables enter a certain relation.

This position is not eliminativist. It does not aim to dissolve consciousness into a description of physical quantities or neural activity. Nor is it dualist. It does not add consciousness as a special entity. It describes consciousness as the dynamic relation among multiple processing layers.

\textbf{Third: redescription of conditions, not a negation of neuroscience}

This paper holds that consciousness does not localize to a single site. This is not a negation of neuroscience. It is entirely possible that specific brain regions, circuits, oscillations, or activity patterns correlate with conscious states. In the framework of this paper, these are interpreted as local conditions that participate in the log-referencing loop.

Neural correlates of consciousness are not consciousness itself. They are part of the local substrate through which consciousness opens. The shift from "finding the seat of consciousness" to "investigating the conditions under which consciousness opens" is the stance this paper takes toward neuroscience.

The reader is invited to hold these three points throughout the discussion that follows. What this paper claims is not the negation or mystification of consciousness, but the fixing of a structural position from which consciousness can be questioned.

Note also that the two-layer structure of fast and slow layers treated in this paper stands on IFGT (Information Field Geometry Theory), Bio-IFGT, and the cognitive-window projection developed in Conceptual Topography. IFGT supplies the general language of information density, flow, and potential; Bio-IFGT explains why such layered dynamics can appear in biological systems; CT describes how conceptual terrains stabilize within a cognitive observation window. The grounds for why this two-layer structure appears in biological systems are given in §4. Even for readers approaching this paper as a standalone work, knowing this connection deepens the meaning of the equations.
